\documentclass[twocolumn]{aastex701}
\begin{document}
\title{A Fictional Hot Jupiter with a Puffed-up Atmosphere Around an Old Metal-poor Star}
\author[orcid=0000-0000-0000-0001]{Hinata Kumoi}
\affiliation{Astronomical Institute, Kitamachi University, Kitamachi 100-0000, Japan}
\email[show]{kumoi@example.org}
\author{Rafael Montcada}
\affiliation{Observatorio Imaginario, Camino del Cielo s/n, Ciudad Ficticia, Freedonia}
\email{montcada@example.org}
\author{Ingrid Halvorsen}
\affiliation{Institute of Fictional Exoplanets, University of Nowhere, Nowhere 00000, Freedonia}
\email{halvorsen@example.org}
\begin{abstract}
We report the discovery of KU-7b, a fictional hot Jupiter transiting an old, metal-poor K dwarf every 3.2 days.
Combining transit photometry with radial velocities, we measure a mass of $0.41\,M_{\rm J}$ and a radius of $1.52\,R_{\rm J}$, making KU-7b one of the least dense giant planets known around a metal-poor host.
The inflated radius cannot be explained by stellar irradiation alone and suggests additional heating in the planet's interior.
Low-resolution transmission spectroscopy reveals a strong sodium feature and a featureless continuum at longer wavelengths, consistent with high-altitude clouds.
KU-7b is an excellent target for atmospheric characterization with current space telescopes.
\end{abstract}
\keywords{\uat{Exoplanets}{498} --- \uat{Hot Jupiters}{753} --- \uat{Transit photometry}{1709}}
\section{Introduction}
Hot Jupiters are giant planets on orbits of a few days, and many of them have radii larger than predicted by models of irradiated planets \citep{kumoi2021,montcada2022}.
The occurrence of such planets correlates strongly with the metallicity of the host star, so giant planets around metal-poor stars are rare and valuable.
They test whether the mechanisms that inflate planetary radii depend on the composition of the planet \citep{halvorsen2023}.

In this Letter we present KU-7b, a fictional inflated hot Jupiter orbiting a metal-poor K dwarf.
Section~2 describes the discovery photometry and the radial-velocity follow-up, and Section~3 presents the transmission spectrum.

\section{Observations}
The transits were first detected in a fictional wide-field survey.
\bibliography{refs}{}
\bibliographystyle{aasjournalv7}
\end{document}
